\markright{\mititulo \hspace{5cm} \thepage}
\pagestyle{fancy}

\section*{Resumen}
\addcontentsline{toc}{section}{Resumen}

\noindent Se identifican los principales círculos y puntos de referencia de la esfera
celeste y se registra la posición del Sol en los sistemas de coordenadas horizontal y
ecuatorial en la plazoleta central de la Universidad de Antioquia, Medellín, Colombia
($6{,}267^\circ$N, $75{,}569^\circ$O, 1495~m.s.n.m.).
La práctica consta de dos partes.
En la \textbf{Parte~1} se usó la herramienta \textit{Solar School House Solar Site
Analysis} \citep{SolarSchoolHouse2026} para registrar cuatro sesiones vespertinas entre
el 3 y el 8 de abril de 2026; los datos son compatibles con las efemérides de JPL
Horizons \citep{Giorgini1996} dentro de $\leq 2\sigma$ en ambos sistemas de coordenadas.
En la \textbf{Parte~2} se midió la longitud de sombra proyectada por un gnomon de
$L = (102{,}0 \pm 0{,}1)$~cm al mediodía civil (12:00 COL) durante seis días entre el
10 y el 18 de abril de 2026; a partir de esas sombras se derivaron la elevación solar y
la declinación del Sol mediante las expresiones trigonométricas del gnomon.
Las declinaciones derivadas concuerdan con los valores de referencia de JPL
Horizons \citep{Giorgini1996} dentro de la incertidumbre instrumental
($\sigma_\delta = 0{,}45^\circ$).

\vspace{1cm}\textit{Palabras clave:} esfera celeste, coordenadas celestes, sistema horizontal, sistema ecuatorial, gnomon, declinación, culminación solar

\newpage

\section*{Abstract}
\addcontentsline{toc}{section}{Abstract}

\noindent The main circles and reference points of the celestial sphere are identified,
and the Sun's position is recorded in the horizontal and equatorial coordinate systems
at the central plaza of the Universidad de Antioquia, Medellín, Colombia
($6.267^\circ$N, $75.569^\circ$W, 1495~m.a.s.l.).
The practice has two parts.
In \textbf{Part~1}, the \textit{Solar School House Solar Site Analysis} tool
\citep{SolarSchoolHouse2026} was used to observe four afternoon sessions between April~3
and~8, 2026; the data are compatible with JPL Horizons ephemerides \citep{Giorgini1996}
within $\leq 2\sigma$ in both coordinate systems.
In \textbf{Part~2}, the shadow length cast by a gnomon of
$L = (102.0 \pm 0.1)$~cm was measured at civil noon (12:00 COL) on six days between
April~10 and~18, 2026; the solar elevation and declination were then derived from those
shadow lengths using the trigonometric relations of the gnomon.
The derived declinations agree with the JPL Horizons reference values
\citep{Giorgini1996} within the instrumental uncertainty ($\sigma_\delta = 0.45^\circ$).

\vspace{1cm}\textit{Keywords:} celestial sphere, celestial coordinates, horizontal system, equatorial system, gnomon, declination, solar culmination

\newpage
